% TOP_PROBLEM: {"id": 3399, "title": "Complexity of square-root sum", "category_id": 15, "category": {"id": 15, "name": "computer_science", "display_name": "Computer Science", "description": "Computational complexity, algorithms, and theoretical CS.", "slug": "computer-science", "order_index": 15, "created_at": "2026-07-31T15:26:25.671Z"}, "status": "open", "problem_number": "OPG-467", "set_id": 12}
% FIRST_POSTED: 2026-10-01
% ENTRY_KIND: statement_only
% UnsolvedMath ID 3399; source catalogue code OPG-467
% !TeX program = lualatex
% Definition and sources only; this file is not an LLM solution attempt.
\documentclass[11pt,a4paper]{article}
\usepackage[margin=25mm]{geometry}
\usepackage{fontspec}
\setmainfont{lmroman10-regular.otf}[BoldFont=lmroman10-bold.otf,
 ItalicFont=lmroman10-italic.otf,BoldItalicFont=lmroman10-bolditalic.otf]
\usepackage{amsmath,amssymb,mathtools}
\usepackage{unicode-math}
\setmathfont{latinmodern-math.otf}
\AtBeginDocument{\let\setminus\smallsetminus}
\usepackage{xurl,hyperref}
\hypersetup{colorlinks=true,urlcolor=blue}
\setcounter{secnumdepth}{0}
\setlength{\parindent}{0pt}
\setlength{\parskip}{5pt}
\setlength{\emergencystretch}{3em}
\begin{document}
\section*{Complexity of square-root sum}\label{3399}
{\small UnsolvedMath ID 3399 · OPG-467 · Computer Science · OpenGarden}
\subsection{Definitions and mathematical statement}
The square-root-sum decision problem takes nonnegative integers
$a_1,\ldots,a_n$ and a nonnegative integer threshold $k$, all encoded
in binary. It asks whether the real-number inequality
\[
                            \sum_{i=1}^n\sqrt{a_i}\le k
\]
holds, where each square root is the nonnegative one. Input size is
the total number of bits in the list and threshold. Equality is
included among yes-instances. The problem is to determine the
classical worst-case computational complexity of this language,
in particular whether it lies in $\mathrm P$ or at least in
$\mathrm{NP}$.

Membership in $\mathrm P$ would mean a deterministic Turing machine
deciding all instances in time polynomial in the input length.
Membership in $\mathrm{NP}$ would mean that every yes-instance has
a polynomial-length certificate verifiable in polynomial time, with
no such certificate accepted on a no-instance.
These are exact bit-complexity questions: arithmetic with arbitrary
real numbers is not a unit-cost operation. Numerical approximations
are useful only when accompanied by a separation or equality argument
that guarantees the correct comparison within the claimed time.
The number of summands varies with the input and is not fixed.
Bounds polynomial in the values $a_i$ rather than their bit lengths
do not constitute polynomial-time algorithms in this formulation.
\subsection{Short English statement}
How hard is it to decide exactly whether a sum of square roots of binary-encoded nonnegative integers is at most a given integer?
\subsection{Sources}
\begin{itemize}
\item[S1] ULAM.AI and the UnsolvedMath Contributors, supplied problems.json, record 3399 (OPG-467), OpenGarden collection. Catalogue identity and source transcription; CC BY 4.0.
\url{https://huggingface.co/datasets/ulamai/UnsolvedMath}.
\item[S2] Open Problem Garden, Complexity of square-root sum. Original problem and discussion; the mathematical formulation below is expanded from the supplied source record.
\url{http://www.openproblemgarden.org/op/complexity_of_square_root_sum}.
\end{itemize}
\end{document}
